\chapter{Manufacturing Architecture}
\label{ch:architecture}

\section{Fortune-5 FIBO LineController Factory}

The experimental apparatus consists of four layers:

\subsection{Layer 1: FIBO-Aligned Control Ontology}

The base ontology (\texttt{f5\_line\_control.ttl}) defines:
\begin{itemize}
    \item 30+ connector types (CRM, KYC/AML, Credit Bureau, Treasury, etc.)
    \item 20 operations per connector = 600 operation modules
    \item Generation patterns specifying LOC multipliers
    \item FIBO references for semantic grounding
\end{itemize}

\begin{lstlisting}[language=]
f5:CRMConnector a ln:Connector ;
    ln:connectorId "crm" ;
    ln:operations ( f5:CreateLeadOp f5:UpdateLeadOp ... ) ;
    ln:authScheme "oauth2" ;
    ln:rateLimit "1000/min" .
\end{lstlisting}

\subsection{Layer 2: SPARQL Extraction}

SPARQL queries traverse the ontology graph, extracting connector metadata:

\begin{lstlisting}[language=SQL]
SELECT ?connectorId ?name ?authScheme ?rateLimit
       (GROUP_CONCAT(?operation; separator=",") AS ?operations)
WHERE {
    ?connector a ln:Connector ;
               ln:connectorId ?connectorId ;
               ln:operations ?opList .
    ?opList rdf:rest*/rdf:first ?operation .
}
\end{lstlisting}

This computes the parallel transport along semantic edges.

\subsection{Layer 3: Tera Template Engine}

Templates define the fiber structure---how concepts map to code:

\begin{lstlisting}
-module(f5_connector_{{ connectorId }}).
{% for op in operations %}
{{ op | lower }}(Params) ->
    gen_server:call(?MODULE, { {{op | lower }}, Params}).
{% endfor %}
\end{lstlisting}

\subsection{Layer 4: OTP Evidence Harness}

The \texttt{f5\_evidence} module collects runtime proofs using OTP-native calls:

\begin{itemize}
    \item \texttt{erlang:system\_info/*} for VM state
    \item \texttt{erlang:process\_info/*} for scheduler statistics
    \item \texttt{crypto:hash/2} for cryptographic chaining
\end{itemize}

This ensures evidence cannot be faked---it requires actual BEAM VM execution.
